% Part: first-order-logic
% Chapter: sequent-calculus
% Section: rules-and-proofs

\documentclass[../../../include/open-logic-section]{subfiles}

\begin{document}

\iftag{FOL}
      {\olfileid{fol}{seq}{rul}}
      {\olfileid{pl}{seq}{rul}}

\olsection{قاعدې او \usetoken{P}{derivation}}

د راتلونکو خبرو دپاره، پرېږدئ چې $\Gamma, \Delta, \Pi, \Lambda$ د
!!{sentence}s متناهي لړۍ وښيي۔

\begin{defn}[سېکوېنټ]
\emph{سېکوېنټ} د لاندې بڼې څرګندنه ده:
\[
\Gamma \Sequent \Delta
\]
دلته $\Gamma$ او $\Delta$ د ژبې $\Lang L$ د !!{sentence}s متناهي
(او کېدے شي تشې) لړۍ دي۔ $\Gamma$ ته \emph{مقدم} او $\Delta$ ته
\emph{تالي} ويل کېږي۔
\end{defn}

\begin{explain}
د يوه سېکوېنټ تر شا شهودي مفکوره دا ده: که د مقدم ټولې !!{sentence}s
صدق کوي، نو د تالي لږ تر لږه يوه !!{sentence}s صدق کوي۔ يعنې که
$\Gamma = \tuple{!A_1, \dots, !A_m}$ او
$\Delta = \tuple{!B_1, \dots, !B_n}$ وي، نو $\Gamma \Sequent \Delta$
هله صدق کوي چې
\[
(!A_1 \land \cdots \land !A_m) \lif (!B_1 \lor \cdots \lor
!B_n)
\]
صدق کوي۔ دوه خاص حالتونه شته: يو هغه چې $\Gamma$ تش وي، او بل هغه
چې $\Delta$~تش وي۔ کله چې $\Gamma$ تش وي، يعنې $m = 0$، نو $\quad
\Sequent \Delta$ هله صدق کوي چې $!B_1 \lor \dots \lor !B_n$ صدق
کوي۔ کله چې $\Delta$ تش وي، يعنې $n = 0$، نو $\Gamma \Sequent \quad$
هله صدق کوي چې $\lnot(!A_1 \land \dots \land !A_m)$ صدق کوي۔
موږ يو سېکوېنټ هله معتبر بولو چې ورسره سمون لرونکې !!{sentence}
معتبره وي۔
\end{explain}

که $\Gamma$ د !!{sentence}s يوه لړۍ وي، نو $\Gamma, !A$ د هغې
پايلې دپاره ليکو چې~$!A$ د~$\Gamma$ ښي پاے ته ورزيات شي (او $!A,
\Gamma$ د هغې پايلې دپاره چې~$!A$ د~$\Gamma$ کيڼ پاے ته ورزيات
شي)۔ که $\Delta$ هم د !!{sentence}s يوه لړۍ وي، نو $\Gamma,
\Delta$ د دواړو لړيو نښلول دي۔

\begin{defn}[لومړنے سېکوېنټ]
\emph{لومړنے سېکوېنټ} داسې سېکوېنټ دے
\iftag{prvFalse,prvTrue}{چې له لاندې بڼو څخه يوه لري:
  \begin{enumerate}
    \item $!A \Sequent !A$
    \tagitem{prvTrue}{$\quad \Sequent \ltrue$}{}
    \tagitem{prvFalse}{$\lfalse \Sequent \quad$}{}
  \end{enumerate}}
{چې بڼه ئې $!A \Sequent !A$ ده}، او $!A$ د ژبې هره !!{sentence}
کېدے شي۔
\end{defn}

په سېکوېنټ حساب کښې !!^{derivation}s د سېکوېنټونو داسې ټاکلې ونې
دي چې تر ټولو بره سېکوېنټونه ئې لومړني سېکوېنټونه وي؛ او که يو
سېکوېنټ د يوه يا دوو نورو سېکوېنټونو لاندې راشي، بايد د استنتاج د
يوې قاعدې له مخې ترې په سمه توګه راووځي۔ د~$\Log{LK}$ قاعدې پر دوو
اصلي ډولونو وېشل کېږي: \emph{منطقي} قاعدې او \emph{جوړښتي} قاعدې۔
منطقي قاعدې د لاندې سېکوېنټ د هغې !!{sentence} د !!{main
  operator} په نوم نومول کېږي چې $!A$ او/يا $!B$ لري۔ هره قاعده
دوه بڼې لري: يوه د داسې سېکوېنټ د استنتاج دپاره چې د
!!{operator} لرونکې !!{sentence} ئې کيڼ لوري ته وي، او بله د هغه
سېکوېنټ دپاره چې دا !!{sentence} ئې ښي لوري ته وي۔

\end{document}
